\documentclass{article}

\usepackage{microtype}
\usepackage{graphicx}
\usepackage{subfigure}
\usepackage{booktabs}
\usepackage{amsmath}
\usepackage{amssymb}
\usepackage{hyperref}

\newcommand{\theHalgorithm}{\arabic{algorithm}}

\usepackage[accepted]{icml2021}

\icmltitlerunning{Fewer Fruits, Fewer Bites: SARSA vs Q-learning on Snake}

\begin{document}

\twocolumn[
\icmltitle{Fewer Fruits, Fewer Bites:\\SARSA vs Q-learning on Snake under Full and Partial Observability}

\begin{icmlauthorlist}
\icmlauthor{Matteo Polo}{}
\end{icmlauthorlist}

\icmlkeywords{Reinforcement Learning, SARSA, Q-learning, Reward Design, Partial Observability, UniPD}

\vskip 0.3in
]

%\printAffiliationsAndNotice{}

\begin{abstract}
We train a convolutional Q-network to play Snake with SARSA and Q-learning.
Our first implementation ate about 72 fruits per episode, but lost almost all of them in about 6 self-collisions.
We therefore shifted the goal of the project to eating as many fruits as possible \emph{while avoiding self-collisions} as much as possible.
% and redesigned the reward, the allowed moves and the training loop accordingly.
Compared with the first implementation, self-collisions drop by about 70\% and lost segments by about 75\%, at the cost of about 35\% fewer fruits.
We then compare the two algorithms with the full board and with a $5\times5$ window: a cosine learning-rate decay removes most of the instability of the greedy policy, after which all four configurations eat about 40 fruits per game by the end of training, while Q-learning takes more risk than SARSA during training.
\end{abstract}

% ---------------------------------------------------------------------------
\section{Introduction}
\label{sec:intro}

Snake is a small control problem with a delayed, competing objective: eating a fruit makes the snake longer, and a longer snake is harder to steer without biting itself.
Our project had two goals: (i) eat as many fruits as possible, (ii) avoid self-collisions; and one question: how do on-policy SARSA \cite{rummery1994} and off-policy Q-learning \cite{watkins1992} differ on this trade-off when the agent sees the whole board versus only a window around its head?

Most of the work went into \emph{making the agents learn the intended behaviour} rather than into the algorithms themselves.
Each fix exposed a new failure mode, so the paper is organised as a sequence of design decisions (Section~\ref{sec:design}), each motivated by an observed failure.
Our contributions are:
\begin{itemize}
\item a diagnosis of reward hacking in the provided reward and a corrected collision penalty that makes self-collisions strictly unprofitable (Section~\ref{sec:hacking});
\item a set of minimal changes (action masks, step cost, per-episode reset) that turn a degenerate policy into one that grows the snake (Sections~\ref{sec:collapse}--\ref{sec:reset});
\item evidence that the instability of the greedy policy between checkpoints comes from the fixed learning rate, and that a cosine decay removes most of it (Section~\ref{sec:lr});
\item a controlled $2\times2$ comparison (algorithm $\times$ observability) in which the two scripts differ \emph{only} in the TD target and the two environments differ \emph{only} in the observation, with a rule-based baseline as reference (Section~\ref{sec:results}).
\end{itemize}

% ---------------------------------------------------------------------------
\section{Setup}
\label{sec:setup}

\subsection{Environment}
Each board is a $10\times10$ grid surrounded by walls, i.e.\ 64 free cells (Figure \ref{fig:map}).
At every step the agent picks one of four moves.
Three facts about the setup matter for everything that follows:
\begin{itemize}
\item a \textbf{self-collision does not end the game}: the tail is cut from the bitten segment onwards and the snake keeps moving;
\item a \textbf{wall hit does not end the game} either: the head stays still while the tail retracts by one cell;
\item 32 boards are simulated in parallel, and an \emph{episode} is a window of 500 steps.
\end{itemize}


\begin{figure}[htp]
    \centering
    \includegraphics[width=0.5\columnwidth]{figures/snake_map.jpg}
    \caption{A $10\times10$ board: walls (black), head (light green), body (dark green), fruit (red).}
    \label{fig:map}
\end{figure}

Since nothing is terminal, the end of an episode is a time-limit truncation, and we keep bootstrapping on the last transition \cite{pardo2018}.

\subsection{State representation}
The board is one-hot encoded into four channels (empty, fruit, body, head); walls are the all-zero vector.
\begin{itemize}
\item \textbf{Fully observable (FO):} the whole $10\times10\times4$ board.
\item \textbf{Partially observable (PO):} a $(2m+1)\times(2m+1)\times4$ window centred on the head, with $m=2$ ($5\times5$). Cells outside the board are padded as walls.
\end{itemize}
%The two environment files are identical except for this observation: with the same seed and actions they produce the same boards and rewards (checked over 3000 steps on 32 boards).
%We chose $m=2$ because it always contains the cells the agent can reach in two moves, so immediate dangers are visible, while the fruit is usually not.

\subsection{Original implementation}
Our first implementation, which we call the \emph{original} agents, used the environment reward as given (Table~\ref{tab:reward}), no action masks (reversals and wall hits allowed), games that continued across episodes, and a fixed learning rate. Its two scripts also used different $\varepsilon$ schedules (SARSA $0.2\to0.02$, Q-learning constant $0.15$). Network, heuristic, 32 boards and 200 episodes of 500 steps are the same as in the modified agents. Section~\ref{sec:design} motivates every change, and Section~\ref{sec:results} measures the result.

\subsection{Heuristic baseline}
\label{sec:baseline}
As a reference without any learning we use a greedy rule-based policy. At every step it looks at the four cells next to the head and:
\begin{enumerate}
\item discards the reversal and the moves into a wall (if no move is left, it reverses);
\item among the remaining moves, prefers those that do not enter the body (if all do, it keeps them all);
\item among these, picks the one that minimises the Manhattan distance between head and fruit, breaking ties at random.
\end{enumerate}
In PO it sees only the $5\times5$ window: if the fruit is outside, it skips rule 3 and picks a random safe move.
This is the natural ``go to the fruit without crashing'' strategy: it never hits a wall and never bites itself when a free cell is available, but it looks only one step ahead, so a long snake can enter a region closed by its own body.
We evaluate it with the same protocol as the greedy evaluation of the agents (Section~\ref{sec:reset}): 20 new games of 500 steps with a fixed seed.
The same rules also drive the exploration of the agents (Section~\ref{sec:netexp}), so the comparison asks whether the network learns more than the rule it explores with.

\subsection{Network and exploration}
\label{sec:netexp}
Both algorithms share the same Q-network: two $3\times3$ convolutions (32 and 64 filters, \emph{same} padding), a dense layer of 128 units and a linear output with 4 Q-values.
In PO only the input size changes.
We train online with Adam, with a learning rate that decays from $10^{-3}$ to $10^{-4}$ along a cosine (Section~\ref{sec:lr}), MSE on the TD error, one gradient step per environment step on the 32 transitions, $\gamma=0.95$, for 200 episodes.

Exploration is $\varepsilon$-greedy with $\varepsilon$ decaying linearly from $0.2$ to $0.02$, but the exploratory move is not uniform: it is the move of the baseline heuristic (Section~\ref{sec:baseline}), which in PO also sees only the window.
Uniform exploration on a crowded board mostly produces collisions, which would drown the learning signal.

% ---------------------------------------------------------------------------
\section{Algorithms}
\label{sec:algo}

The two methods differ only in the TD target.
Let $\mathcal{A}(s)$ be the set of \emph{allowed} actions in $s$ (no wall, no reversal; Section~\ref{sec:collapse}). Then
\begin{align}
y^{\textsc{sarsa}} &= r + \gamma\, Q(s',a'), \quad a' \sim \pi_\varepsilon(\cdot \mid s'), \label{eq:sarsa}\\
y^{\textsc{ql}} &= r + \gamma \max_{a \in \mathcal{A}(s')} Q(s',a). \label{eq:ql}
\end{align}
In both cases $a'$ is chosen only among the allowed actions $\mathcal{A}(s')$. SARSA evaluates the behaviour policy, exploration included. Q-learning evaluates the greedy policy.
The $\max$ in \eqref{eq:ql} is restricted to $\mathcal{A}(s')$; otherwise the target would bootstrap on moves the agent can never take.
The training loop is summarised in Algorithm~\ref{alg:train}.
%the two scripts are generated from the same code, and a \texttt{diff} confirms that only the target computation and file names differ.

\begin{algorithm}[tb]
   \caption{Training loop (shared by SARSA and Q-learning)}
   \label{alg:train}
\begin{algorithmic}
   \FOR{episode $=1$ {\bfseries to} $200$}
   \STATE Reset the 32 boards; set $\varepsilon$ (linear decay)
   \STATE Choose $a$: greedy over $\mathcal{A}(s)$, heuristic w.p.\ $\varepsilon$
   \FOR{$t=1$ {\bfseries to} $500$}
   \STATE Execute $a$; observe environment reward and $s'$
   \STATE $r \leftarrow$ reward \eqref{eq:collision} $+$ step cost
   \STATE Choose $a'$ in $s'$ as above
   \STATE $y \leftarrow$ \eqref{eq:sarsa} (SARSA) or \eqref{eq:ql} (Q-learning)
   \STATE One Adam step on $\frac{1}{32}\sum_{b}\big(y_b - Q(s_b,a_b)\big)^2$
   \STATE $s \leftarrow s'$, $a \leftarrow a'$
   \ENDFOR
   \STATE Every 10 episodes: greedy evaluation on 20 new games
   \ENDFOR
\end{algorithmic}
\end{algorithm}

% ---------------------------------------------------------------------------
\section{Design decisions}
\label{sec:design}

Table~\ref{tab:reward} summarises the final reward. Each subsection explains one change through the failure that motivated it, in the order in which we made them. Sections~\ref{sec:hacking}--\ref{sec:reset} were tested with SARSA in FO on short 60-episode runs, and their numbers come from those runs, except those of the original agents; the learning-rate decay (Section~\ref{sec:lr}) was tested on all four configurations with 200 episodes.

\begin{table}[t]
\caption{Reward per step. $k$ = segments lost in a self-collision.}
\label{tab:reward}
\vskip 0.1in
\begin{center}
\begin{small}
\begin{tabular}{lcc}
\toprule
Event & Original & Modified \\
\midrule
Fruit eaten        & $+0.5$      & $+0.5$ \\
Self-collision     & $-0.1\,k$   & $-(0.5+0.5\,k)$ \\
Wall hit           & $-0.3$      & masked \\
Any step           & $0$         & $-0.01$ \\
\bottomrule
\end{tabular}
\end{small}
\end{center}
\vskip -0.15in
\end{table}

\subsection{Reward hacking: ``eat, then bite''}
\label{sec:hacking}
With the provided reward the original agents reached a stable, positive return (about $+29$ per episode), yet the snake stayed short: per episode they ate about 72 fruits and lost about 71 segments in about 6 self-collisions (Figure~\ref{fig:origmod}).
The reason is arithmetic. A segment is worth $+0.5$ when gained but costs only $-0.1$ when lost, so the cycle ``eat $k$ fruits, then cut $k$ segments'' yields
\begin{equation}
0.5k - 0.1k = +0.4k > 0 .
\end{equation}
Discounting makes it worse, because the fruit comes before the penalty, and a short snake also reaches fruits faster.
This is specification gaming \cite{amodei2016,krakovna2020}: the agent maximises the reward it is given, and that reward disagrees with goal (ii).
Potential-based shaping cannot fix this, since by construction it preserves the optimal policy \cite{ng1999}; the reward itself has to change.
We replace the collision penalty with
\begin{equation}
r_{\text{self}}(k) = -(0.5 + 0.5\,k).
\label{eq:collision}
\end{equation}
The proportional part gives back exactly what the lost segments were worth, so the cycle becomes neutral; the fixed part makes it strictly negative, $0.5k-(0.5+0.5k)=-0.5$.
When a collision is unavoidable, the dependence on $k$ still prefers biting close to the tail.

%We also found that our original implementation counted self-collisions from the reward value, so a collision costing exactly 3 segments ($3\times(-0.1)=-0.3$) was mistaken for a wall hit: 12.6\% of collisions in a test run.
%All our metrics are computed from the board state instead, so \emph{the original collision counts are underestimates}.

\subsection{Passive collapse: action masks and step cost}
\label{sec:collapse}
With \eqref{eq:collision} alone, the agent stopped eating: a snake of length 1 oscillated between two cells for the whole episode, with return 0.
At length 1 a reversal is free; after one fruit the same reversal bites the neck, worth $0.5-1.0=-0.5$. Not eating (value 0) is therefore a local optimum of the learning dynamics, not of the task.
Two changes remove it:
\begin{itemize}
\item \textbf{Reversal mask.} The move onto the neck is impossible in real Snake, so we set its Q-value to $-\infty$ before the $\arg\max$. At length 1 we mask the opposite of the last direction.
\item \textbf{Step cost} $-0.01$. Wandering forever is now worth $-0.01/(1-\gamma)=-0.2$ instead of 0.
\end{itemize}
Bites on the rest of the body are \emph{not} masked: avoiding traps is what the network must learn.

\subsection{Wall loops}
\label{sec:walls}
After the previous fix the snake grew (mean length from 3.5 to 35.6), but the training return \emph{decreased}. Splitting the return into its terms showed why: wall hits rose from about 5 to 45 per episode.
A wall hit leaves the state almost unchanged (the head does not move), so a greedy choice of ``wall'' is repeated at the next step; at length 1 the loop is infinite.
We mask wall moves too \cite{huang2022}, so that $\mathcal{A}(s)$ excludes walls and reversals, consistently with the heuristic.
Wall hits dropped to 0. At steady state the return is $\approx -5 - 0.5\cdot(\text{collisions})$, because fruits and lost segments cancel; the training return is therefore a poor progress measure, and we report task metrics instead.

\subsection{Distribution shift: per-episode reset and greedy evaluation}
\label{sec:reset}
In our original implementation the games continued across episodes, so after a few episodes the network almost only sees long snakes.
A separate greedy evaluation ($\varepsilon=0$, 20 new games of 500 steps, fixed seed) revealed the problem: the network that ate about 30 fruits per episode in training mostly ate 0--1 fruits in new games, looping with a short snake.
Resetting the boards at every episode keeps the start of the game in the training data. Greedy evaluation went from 9 to 23 fruits at episode 59 of the same configuration, and the per-episode return became comparable across episodes.
Because averages hide bimodal outcomes, the evaluation also counts \emph{stuck} games: games with no fruit in the last 100 steps, typically a loop.

\subsection{Instability: learning-rate decay}
\label{sec:lr}
After the reset, the greedy policy still entered long cycles (period 12--20) after a few fruits, while in training it ate steadily.
In training a few heuristic moves per game (at $\varepsilon\approx0.02$, about 10) are enough to break a cycle, so the loops only appear when the agent acts greedily.
In the 200-episode runs with a fixed learning rate of $10^{-3}$, the greedy policy also alternated good and broken checkpoints a few episodes apart, up to the end of training (grey lines in Figure~\ref{fig:lr}), while the training curves stayed smooth.
Across the checkpoints from episode 100 on, greedy fruits had a standard deviation of 9--15, and 12 of the 44 checkpoints of the four configurations ate fewer than 10 fruits, with most games stuck (Table~\ref{tab:lr}).
The problem is worst in PO, where 6 of 11 Q-learning checkpoints are broken.
Updates are online on correlated transitions, and with a constant step size even late updates can flip the $\arg\max$ in many states at once.
We therefore let the learning rate decay from $10^{-3}$ to $10^{-4}$ along a cosine over the $10^5$ updates of training: early learning is almost unchanged, while late updates are small.
With the decay only one checkpoint out of 44 is broken (Q-learning PO at episode 140), and from episode 150 on every configuration eats 40--43 fruits per game with a standard deviation below 3 and at most 1.3 stuck games of 20, against 18--31 fruits and 7--12 stuck games with the fixed rate (Table~\ref{tab:lr}, Figure~\ref{fig:lr}).
The effect is largest in PO, where greedy fruits roughly double, and training also improves slightly (2--3 more fruits per episode).

\begin{table}[t]
\caption{Greedy evaluation with a fixed vs cosine learning rate (20 games per checkpoint). Broken = checkpoints with fewer than 10 fruits; stuck = games with no fruit in the last 100 steps (of 20); std = standard deviation of fruits across checkpoints.}
\label{tab:lr}
\vskip 0.1in
\begin{center}
\begin{small}
\setlength{\tabcolsep}{3pt}
\begin{tabular}{llccccc}
\toprule
 & & \multicolumn{3}{c}{Checkpoints $\ge100$} & \multicolumn{2}{c}{$\ge150$} \\
\cmidrule(lr){3-5}\cmidrule(lr){6-7}
 & LR & Fruit & Broken & Stuck & Fruit & Std \\
\midrule
SARSA, FO    & fixed  & 29.1 & 0/11 & 7.8  & 30.6 & 10.0 \\
             & cosine & 37.7 & 0/11 & 4.1  & 43.2 & 0.6 \\
Q-learn., FO & fixed  & 21.8 & 3/11 & 10.5 & 26.6 & 13.0 \\
             & cosine & 38.5 & 0/11 & 3.0  & 42.4 & 2.7 \\
SARSA, PO    & fixed  & 23.4 & 3/11 & 9.6  & 24.5 & 13.4 \\
             & cosine & 39.7 & 0/11 & 1.5  & 41.5 & 0.7 \\
Q-learn., PO & fixed  & 14.4 & 6/11 & 14.1 & 17.5 & 14.8 \\
             & cosine & 32.3 & 1/11 & 5.1  & 40.3 & 1.2 \\
\bottomrule
\end{tabular}
\end{small}
\end{center}
\vskip -0.15in
\end{table}

\begin{figure}[t]
\begin{center}
\centerline{\includegraphics[width=\columnwidth]{figures/lr_fixed_vs_cosine.pdf}}
\caption{Greedy evaluation (fruits per game, 20 games) at every checkpoint, with a fixed learning rate (grey) and a cosine decay (colour), one panel per configuration. With the decay the policy stops oscillating as the learning rate becomes small (from about episode 150); with the fixed rate it keeps oscillating until the end.}
\label{fig:lr}
\end{center}
\vskip -0.3in
\end{figure}

% ---------------------------------------------------------------------------
\section{Results}
\label{sec:results}

All runs use the final configuration (learning-rate decay included), 200 episodes and one seed. Training metrics are averaged over the last 50 episodes. Greedy metrics are averaged over the checkpoints from episode 100 on.

\begin{figure*}[t]
\begin{center}
\centerline{\includegraphics[width=\textwidth]{figures/orig_vs_mod.pdf}}
\caption{Original vs modified agents (FO), mean $\pm$ std over the last 50 training episodes. Hatched: original scripts. The modified agents eat about 35\% fewer fruits, but make about 70\% fewer self-collisions and lose about 75\% fewer segments.}
\label{fig:origmod}
\end{center}
\vskip -0.3in
\end{figure*}

\subsection{Original vs modified}
Figure~\ref{fig:origmod} compares the original scripts with the modified ones.
Fruits drop from 72--74 to 46--48 per episode, self-collisions from 5.9--6.6 to 1.8--2.2, and lost segments from 71--72 to 14--19.
The original agents lost almost every segment they gained; ours keep them, and the mean length grows to 22--23.
Measured as fruits per self-collision, efficiency rises from about 12 to 22--26.

The comparison is not perfectly controlled: the original agents do not reset the games between episodes, so an episode measures a different slice of play, they use different $\varepsilon$ schedules and a fixed learning rate (Section~\ref{sec:setup}).
The fruit gap is therefore not only a cost of caution. The effect on collisions and lost segments, however, is large and consistent in both algorithms.

\begin{figure}[t]
\begin{center}
\centerline{\includegraphics[width=\columnwidth]{figures/training_growth.pdf}}
\caption{Training: fruits per episode (top) and mean snake length (bottom), moving average over 10 episodes. Colour: algorithm; solid: fully observable (FO); dashed: partially observable (PO, $5\times5$ window). Black horizontal lines: heuristic baseline (Section~\ref{sec:baseline}).}
\label{fig:growth}
\end{center}
\vskip -0.3in
\end{figure}

\begin{figure}[t]
\begin{center}
\centerline{\includegraphics[width=\columnwidth]{figures/training_safety.pdf}}
\caption{Training: self-collisions (top) and segments lost (bottom) per episode, same encoding as Figure~\ref{fig:growth}, black lines = baseline. Both decrease steadily in the second half of training.}
\label{fig:safety}
\end{center}
\vskip -0.3in
\end{figure}

\begin{table}[t]
\caption{$2\times2$ comparison. Training: mean of the last 50 episodes. Greedy: mean over checkpoints $\ge 100$ (20 games each); Bite = self-collisions, Lost = segments lost, Len.\ = mean snake length (per episode); stuck = games with no fruit in the last 100 steps. The baseline does not learn: its row is one evaluation (20 games of 500 steps), valid for both groups of columns.}
\label{tab:2x2}
\vskip 0.1in
\begin{center}
\begin{small}
\setlength{\tabcolsep}{3.5pt}
\begin{tabular}{lcccccc}
\toprule
 & \multicolumn{4}{c}{Training} & \multicolumn{2}{c}{Greedy} \\
\cmidrule(lr){2-5}\cmidrule(lr){6-7}
 & Fruit & Bite & Lost & Len. & Fruit & Stuck \\
\midrule
SARSA, FO      & 45.9 & 1.77 & 13.7 & 22.9 & 37.7 & 4.1 \\
Q-learn., FO   & 47.5 & 2.15 & 19.3 & 21.8 & 38.5 & 3.0 \\
SARSA, PO      & 42.1 & 1.62 & 19.8 & 18.5 & 39.7 & 1.5 \\
Q-learn., PO   & 41.9 & 1.59 & 19.9 & 18.0 & 32.3 & 5.1 \\
\midrule
Baseline, FO   & 71.8 & 5.90 & 56.0 & 14.0 & 71.8 & 0 \\
Baseline, PO   & 27.6 & 4.05 & 20.2 & 6.4  & 27.6 & 0 \\
\bottomrule
\end{tabular}
\end{small}
\end{center}
\vskip -0.15in
\end{table}

\subsection{SARSA vs Q-learning, FO vs PO}
Table~\ref{tab:2x2} and Figures~\ref{fig:growth}--\ref{fig:eval} report the $2\times2$ design.

\textbf{Fruits vs length.} In FO, fruits per episode peak at about 49--52 between episodes 50 and 80 and then decrease to 46--48 (Figure~\ref{fig:growth}), while the snake keeps growing and collisions keep falling (Figure~\ref{fig:safety}). The agent trades fruits for safety: a longer body needs longer, safer paths to the fruit. The numbers in Table~\ref{tab:2x2} are taken at the end of training, after this trade-off.

\textbf{Algorithm (FO).} Q-learning eats slightly more than SARSA (47.5 vs 45.9) but bites more often (2.15 vs 1.77) and loses about 40\% more segments (19.3 vs 13.7).
This is consistent with the textbook difference \cite{sutton2018}: SARSA learns the value of the exploring policy, so it prices in the occasional wrong move near the body and keeps a safer margin, while Q-learning values the greedy path and takes tighter routes.
In greedy evaluation the two are equivalent (37.7 vs 38.5 fruits).

\textbf{Observability.} During training PO costs little: about 10\% fewer fruits (42 vs 46--48) and a shorter snake (18 vs 22--23). PO collisions are lower partly \emph{because} the snake is shorter.
In greedy evaluation there is no clear gap: SARSA PO averages 39.7 fruits, and Q-learning PO 32.3 only because of one broken checkpoint (episode 140, 2 fruits); from episode 150 on all four configurations eat 40--43 fruits with at most 1.3 stuck games of 20.
%With a $5\times5$ window, many different states produce the same observation whenever the fruit is out of view (perceptual aliasing \cite{kaelbling1998}), and a deterministic memoryless policy can then repeat the same cycle forever \cite{singh1994}.
The decay removes most stuck games also in PO (Table~\ref{tab:lr}), although the reward carries no information about the hidden fruit: in our setting the loops came mainly from unstable updates rather than from the limited view.
One caveat remains: the board is small, so the fruit is often within a few moves of the window.
Under PO the two algorithms are indistinguishable.

\subsection{RL agents vs heuristic baseline}
\label{sec:vsbaseline}
\textbf{FO.} The baseline eats more fruits than the agents (71.8 vs 46--48 per episode in training, 38 in greedy evaluation), but bites itself about 3 times as often (5.90 vs 1.8--2.2) and loses 3--4 times as many segments (56.0 vs 14--19), so its snake stays short (mean length 14.0 vs 22--23; Table~\ref{tab:2x2}, Figures~\ref{fig:growth}--\ref{fig:safety}).
It reproduces, without any learning, the ``eat, then bite'' behaviour of the original agents (Section~\ref{sec:hacking}): with one step of lookahead it walks straight to the fruit and, once long, into regions closed by its own body.
The agents are therefore worse on goal (i) and much better on goal (ii), the trade-off we chose in Section~\ref{sec:intro}.
\textbf{PO.} The agents beat the baseline on every metric: more fruits (42 vs 27.6 in training, 32--40 greedy), less than half the self-collisions, fewer lost segments in greedy evaluation and a snake almost three times as long. When the fruit is out of view the baseline moves at random, while the network has learned where to look for it, even though the reward carries no information about the hidden fruit.
In both settings the agents learn more than the rule they explore with.
Closing the fruit gap in FO would require looking further ahead than one step, for example a larger discount factor, longer training, or giving the agent the free space reachable from the head as an input.

\begin{figure}[t]
\begin{center}
\centerline{\includegraphics[width=\columnwidth]{figures/eval_2x2.pdf}}
\caption{Greedy evaluation ($\varepsilon=0$, 20 new games, fixed seed) of every checkpoint from episode 100 on: one dot per checkpoint, black line = mean. Top: fruits per game; bottom: stuck games. Filled: FO; hollow: PO. Black horizontal lines: baseline fruits (solid FO, dashed PO); the baseline has no stuck games. The outlier at about 2 fruits is the Q-learning PO checkpoint at episode 140.}
\label{fig:eval}
\end{center}
\vskip -0.3in
\end{figure}

% ---------------------------------------------------------------------------
\section{Discussion and limitations}
\label{sec:discussion}

\textbf{Training metrics can mislead.} Each failure in Section~\ref{sec:design} was invisible in at least one standard curve: reward hacking had a rising return, the wall loops had a falling return with improving behaviour, and the distribution shift and the cycles only appeared in greedy evaluation on new games. Logging task metrics (length, lost segments, stuck games) and watching the games was essential.

\textbf{Behaviour vs greedy policy.} With a smart exploration heuristic, part of the competence lives in the exploration, not in the network. Before the learning-rate decay this showed up as a large gap between training and greedy performance, especially in PO; after it, greedy games reach about 40 fruits per 500 steps, close to the 42--48 per episode in training.

\textbf{Instability.} With a fixed learning rate the greedy policy alternated good and broken checkpoints while training curves stayed smooth. Experience replay and a target network \cite{mnih2015} are the standard remedies; we kept the plain online algorithms, to compare SARSA and Q-learning in their textbook form, and a decaying step size was enough. It reduces the instability but does not remove it completely: one checkpoint of Q-learning PO still breaks at episode 140. With a fixed rate, a single checkpoint can therefore look much better or much worse than the agent really is.

\textbf{Limitations.}
\begin{itemize}
\item There is one seed per configuration, so the small SARSA vs Q-learning differences are suggestive, not conclusive.
%\item The original runs undercount collisions and use different $\varepsilon$ schedules.
\item We tried a single learning-rate schedule ($10^{-3}\to10^{-4}$, cosine), with no comparison to other schedules.
\end{itemize}
Natural next steps are several seeds with confidence intervals, a larger board where the window hides more, and giving the PO agent memory (frame stacking or a recurrent network \cite{hausknecht2015}).

% ---------------------------------------------------------------------------
\section{Conclusion}
\label{sec:conclusion}

The provided reward taught both algorithms to eat and bite. Making self-collisions strictly unprofitable, masking moves the game forbids, resetting games and decaying the learning rate produced agents that eat about 35\% fewer fruits but bite about 70\% less and keep their length.
SARSA is the more cautious learner under full observability, as theory predicts.
The instability of the greedy policy came mainly from the fixed learning rate: with a cosine decay all four configurations reach about 40 fruits per game, and a $5\times5$ window costs little on a $10\times10$ board even though the reward says nothing about where the fruit is.

\bibliography{main}
\bibliographystyle{icml2021}

\end{document}
